\section{Visual Encoding: Marks and Channels}

The core of the design space of visual encodings is an \emph{orthogonal}
combination of two aspects: graphical elements called \term{marks}, and
\term{visual channels} that control their appearance. Even complex visual
encodings can be decomposed and analysed in terms of their mark and channel
structure (Munzner, Ch.~5; slides 21--22).

\begin{keybox}[The one-sentence summary]
A \term{mark} is a basic geometric primitive that depicts an item or a link;
a \term{channel} controls the appearance of a mark. The effectiveness of a
channel depends on its \emph{type}: \term{magnitude} (ordered) channels match
ordered data, \term{identity} (categorical) channels match categorical data.
\end{keybox}

\subsection{Marks}

A \term{mark} is a basic graphical element classified by the number of spatial
dimensions it requires:
\begin{itemize}
  \item \term{Point} -- 0D mark. Conveys information only through \emph{position};
        because its area is unconstrained, it \emph{can} be size-coded and
        shape-coded (e.g.\ a bubble in a bubble chart).
  \item \term{Line} -- 1D mark. A bar in a bar chart is a line mark. Its length
        is "taken" by the encoded quantity, so it cannot be size-coded along that
        same direction, but its \emph{width} can carry a second attribute.
  \item \term{Area} -- 2D mark. Both size dimensions are intrinsically fixed by
        its shape, so area marks are typically \emph{not} size- or shape-coded
        (e.g.\ a country on a choropleth map). A 3D \term{volume} mark exists
        but is rarely used.
\end{itemize}

\paragraph{Mark types for items vs.\ links.} For tabular data a mark is an
\term{item}. For network data a mark can be a \term{node} (point/line/area) or a
\term{link}. The two link mark types are: \term{connection} (a line showing a
pairwise relationship) and \term{containment} (an area enclosing/nesting items
to show hierarchy). Links cannot be points; individual items can.

\begin{center}
\begin{tikzpicture}[font=\small, >=Stealth]
  % Top row: item marks -- point, line, area
  % Point (0D)
  \fill[tuwblue] (0.8,1.5) circle (3pt);
  \node[font=\scriptsize, align=center] at (0.8,0.95) {Point\\(0D)};
  % Line (1D) -- a bar
  \fill[tuwblue] (3.0,1.2) rectangle (3.3,1.85);
  \node[font=\scriptsize, align=center] at (3.15,0.95) {Line\\(1D)};
  % Area (2D) -- a filled blob
  \fill[tuwblue] (5.3,1.5) ellipse (10pt and 7pt);
  \node[font=\scriptsize, align=center] at (5.3,0.95) {Area\\(2D)};

  % Bottom row: link marks -- connection, containment
  % Connection: two nodes joined by an edge
  \fill[tuwdark] (0.4,-0.7) circle (2.5pt);
  \fill[tuwdark] (1.8,-0.7) circle (2.5pt);
  \draw[thick, tuwdark] (0.4,-0.7) -- (1.8,-0.7);
  \node[font=\scriptsize, align=center] at (1.1,-1.25) {Connection};
  % Containment: big box with small boxes nested
  \draw[thick, keygreen] (3.6,-1.1) rectangle (5.6,-0.3);
  \fill[keygreen] (3.85,-0.95) rectangle (4.25,-0.45);
  \fill[keygreen] (4.45,-0.95) rectangle (4.85,-0.45);
  \fill[keygreen] (5.05,-0.95) rectangle (5.45,-0.45);
  \node[font=\scriptsize, align=center] at (4.6,-1.5) {Containment};
\end{tikzpicture}
\end{center}
\sketchnote{three rows of icons -- top: a dot (point/0D), a bar (line/1D), a
filled blob (area/2D); bottom: two nodes joined by an edge (connection) and a
big box with small boxes nested inside (containment).}

\subsection{Channels and the magnitude vs.\ identity split}

A \term{visual channel} is a way to control mark appearance, \emph{independent}
of mark dimensionality. The human perceptual system has two fundamentally
different sensory modalities, which splits all channels into two families:

\begin{keybox}[The fundamental split]
\term{Identity / "what--where" channels} (psychophysics: \emph{metathetic})
answer "what is it / where is it?" -- e.g.\ shape, color hue, spatial region,
motion pattern. They suit \term{categorical} attributes. \\[2pt]
\term{Magnitude / "how-much" channels} (psychophysics: \emph{prothetic}) answer
"how much is there?" -- e.g.\ length, area, volume, luminance, saturation, tilt,
position. They suit \term{ordered} (ordinal + quantitative) attributes.
\end{keybox}

\paragraph{List of channels.} Spatial: \term{aligned (common-scale) position},
\term{unaligned position}, \term{depth} (3D position), \term{spatial region}.
Color: \term{hue}, \term{saturation}, \term{luminance}. Size (one per added
dimension): \term{length} (1D), \term{area} (2D), \term{volume} (3D). Motion:
\term{motion pattern}, \term{direction}, \term{velocity}. Plus \term{angle/tilt},
\term{curvature}, and \term{shape} (a complex phenomenon treated as one channel).

\subsection{The two guiding principles}

\begin{keybox}[Expressiveness \& Effectiveness (Mackinlay 1986)]
\textbf{Expressiveness:} the encoding should express \emph{all of, and only,}
the information in the data. Ordered data must be shown so the perceptual system
senses order; unordered data must \emph{not} imply an order that does not exist.
$\Rightarrow$ \emph{match channel type to attribute type}. \\[3pt]
\textbf{Effectiveness:} the \emph{importance} of an attribute should match the
\emph{salience} (noticeability) of its channel. $\Rightarrow$ \emph{most
important attribute $\to$ most effective channel}; less important attributes get
less effective channels.
\end{keybox}

\begin{pitfall}[Classic expressiveness violation]
Encoding an unordered categorical attribute (e.g.\ country) with an ordered
magnitude channel (e.g.\ luminance or length) implies a ranking that does not
exist. This is "a common beginner's mistake." The reverse -- categorical hue
for a quantitative attribute -- destroys perceived order. In both cases the
expressiveness principle is violated even if the chart "works."
\end{pitfall}

\subsection{Channel effectiveness rankings}

Combining the two principles yields two ranked lists (Munzner Fig.~5.6), most
effective at the top. Spatial channels top \emph{both} lists and are the only
ones effective for both data types -- choosing what to encode with position is
the single most central encoding decision, because position-encoded attributes
dominate the user's mental model.

\begin{center}
\begin{tabular}{@{}r l l@{}}
\toprule
Rank & \textbf{Magnitude (ordered)} & \textbf{Identity (categorical)} \\
\midrule
1 & Position on common scale (aligned) & Spatial region \\
2 & Position on unaligned scale        & Color hue \\
3 & Length (1D size)                   & Motion \\
4 & Tilt / angle                       & Shape \\
5 & Area (2D size)                     & \\
6 & Depth (3D position)                & \\
7 & Color luminance                    & \\
8 & Color saturation                   & \\
9 & Curvature                          & \\
10 & Volume (3D size)                  & \\
\bottomrule
\end{tabular}
\end{center}

\noindent Note: luminance and saturation are roughly equally effective, as are
curvature and volume at the bottom.

\subsection{Accuracy: Cleveland \& McGill and Stevens}

\paragraph{Cleveland \& McGill (1984) accuracy ranking.} Controlled experiments
on the magnitude channels gave an empirical ordering of elementary perceptual
tasks, from most to least accurate:
\begin{enumerate}[label=\arabic*.]
  \item Position along a common scale (aligned)
  \item Position along identical, non-aligned scales (unaligned)
  \item Length
  \item Angle / slope
  \item Area
  \item Volume $\approx$ Curvature $\approx$ Luminance (equivalence class)
  \item Color hue \emph{(as a magnitude channel -- last; very different from its
        2nd place as an identity channel)}
\end{enumerate}
Heer \& Bostock (2010) replicated this by crowdsourcing; the one discrepancy is
that \emph{length and angle judgements turned out roughly equivalent}. These
accuracy results are the primary basis for the channel rankings above.

\paragraph{Stevens' power law.} Psychophysics shows that the apparent magnitude
$S$ of a sensory channel follows a power function of the physical intensity $I$:
\[
  S = I^{\,n}.
\]
The exponent $n$ depends on the modality and ranges from sublinear $0.5$
(brightness) to superlinear $3.5$ (electric current).
\begin{itemize}
  \item $n = 1$ (\term{length}): perception is accurate -- doubling the stimulus
        doubles the perception.
  \item $n < 1$ \term{compressive} (e.g.\ \term{area} $\approx 0.7$,
        \term{brightness} $\approx 0.5$): the channel is perceptually
        \emph{compressed} $\Rightarrow$ users \emph{underestimate} large values.
        Doubling brightness looks much less than twice as bright.
  \item $n > 1$ \term{expansive} (e.g.\ red--gray \term{saturation}, electric
        current): the channel is \emph{magnified} $\Rightarrow$ users
        \emph{overestimate}.
\end{itemize}

\begin{keybox}[Why position beats length -- Weber's Law]
Perception is fundamentally \emph{relative}, not absolute (\term{Weber's Law}:
$\delta I / I = K$). Unframed, unaligned bars are hard to compare; adding a
\emph{common frame} or \emph{alignment} gives a common scale, turning a hard
length judgement into an easy position judgement. This is why aligned position
on a common scale ranks first.
\end{keybox}

\subsection{Channel characteristics}

Effectiveness is analysed along five criteria: accuracy, discriminability,
separability, popout, and grouping.

\paragraph{Discriminability.} How many distinguishable \term{bins} (steps) does
a channel provide? The number of attribute values shown must not exceed the
channel's bin count. Example: \term{line width} has very few bins -- great for
3--4 values, poor for hundreds (it starts looking like an area). If they
mismatch, aggregate the attribute or pick another channel.

\paragraph{Separability vs.\ integrality.} Channel pairs lie on a continuum from
\term{separable} (can be attended independently) to \term{integral}
(inextricably fused; the brain perceives an unanticipated combination).

\begin{center}
\begin{tabular}{@{}l l l@{}}
\toprule
Channel pair & Interference & Result \\
\midrule
Position $+$ hue            & none (fully separable) & two clean groupings \\
Size $+$ hue                & some                   & hue hard to read on small marks \\
Width $+$ height            & significant (integral) & fused into perceived \emph{area} (small/med/large) \\
Red $+$ green (RGB axes)    & major (integral)       & perceived as one combined color \\
\bottomrule
\end{tabular}
\end{center}

\noindent Integral is not "bad": horizontal$+$vertical size is a \emph{good}
choice if you want a single attribute with three groups (small, flattened,
large). Match the pair's character to the information.

\paragraph{Popout (preattentive processing).} A distinct item "pops out" from
distractors in time \emph{independent of the number of distractors} -- our
low-level visual system processes the channel in massively parallel fashion
(also called preattentive / tunable detection). Channels supporting popout
include color hue, shape, tilt, size, proximity, motion, shadow direction.
\begin{pitfall}[Popout does not combine]
Popout works for \emph{one channel at a time}. A red circle does \emph{not} pop
out from a field mixing \{red, blue\} $\times$ \{circle, square\} -- it requires
\term{serial search} (time grows linearly with distractors). A few pairs are
exceptions (space$+$color, motion$+$shape); never count on popout for three or
more channels. Also, parallelism is \emph{not} preattentive.
\end{pitfall}

\paragraph{Grouping.} Strength of perceptual grouping cues, strongest first:
\term{containment} (area) $>$ \term{connection} (line) $>$ \term{proximity}
(spatial region) $>$ \term{similarity} (hue, motion, carefully chosen shape).
This proximity strength is exactly why spatial region is the top categorical
channel. Use shape and motion with care: a forward vs.\ backward "C" does not
form a selectable group, while circle vs.\ star does.

\subsection{Channel interference \& redundant encoding}

Two ideas often confused:
\begin{itemize}
  \item \term{Interference} -- encoding \emph{different} attributes in
        non-separable channels (size$+$hue, width$+$height) so the user cannot
        recover each attribute. Avoid.
  \item \term{Redundant encoding} -- encoding \emph{one} attribute in
        \emph{two} channels at once (e.g.\ both hue \emph{and} shape, or both
        position \emph{and} color) to make it very easily perceived. The cost is
        that channels are "used up," so fewer total attributes can be shown.
\end{itemize}

\begin{exambox}[How this is examined]
\textbf{Visualization Critique (open question).} Given a chart, you must:
\begin{itemize}[leftmargin=1.4em]
  \item Decompose it into \emph{marks} and \emph{channels} per attribute.
  \item Check \emph{expressiveness}: is each attribute's channel type matched to
        its data type? Flag categorical-on-magnitude (false order) or vice versa.
  \item Check \emph{effectiveness}: is the most important attribute on the most
        effective channel (position $>$ length $>$ angle $>$ area $>$ color)?
  \item Invoke \emph{Stevens' power law} to argue over/underestimation: area
        ($n\approx0.7$) and brightness ($n\approx0.5$) are underestimated;
        saturation is overestimated; length ($n=1$) is accurate.
  \item Spot \emph{channel interference} (e.g.\ width vs.\ height read as area;
        small marks make hue unreadable) and suggest separable alternatives or
        redundant encoding.
\end{itemize}
\textbf{MCQ on channel properties.} Be ready to: classify a channel as magnitude
vs.\ identity; order channels by effectiveness; recall the Cleveland \& McGill
ranking; identify separable vs.\ integral pairs; state which channels support
popout; and distinguish redundant encoding from interference.
\end{exambox}

\begin{pitfall}[Common confusions to avoid]
Hue is \emph{2nd} as an identity channel but \emph{last} as a magnitude channel.
"Position on a common scale" $\neq$ "length" -- the frame/alignment is what makes
position more accurate (Weber's Law). Effectiveness $=$ importance$\to$salience;
expressiveness $=$ data type$\to$channel type. Do not mix them up.
\end{pitfall}
